% ==============================================================================
% Figura: Diagrama Cinematico Puro del Pendulo Rotacional de Furuta
% Representacion vectorial limpia (estilo libro de texto) entre el marco
% inercial F0 = {O0, X0, Y0, Z0} y la articulacion distal.
% Archivo modular listo para \input{} en IEEEtran (columna simple, figure [!t]).
% Bibliotecas TikZ requeridas: arrows.meta, positioning, fit, calc,
%                            backgrounds, babel, bending.
% ==============================================================================
\begin{figure}[!t]
\centering
% Paleta cromatica de ingenieria (local a la figura)
\definecolor{furutaSlate}{HTML}{1E293B}
\definecolor{furutaLinkEdge}{HTML}{1E4E8C}
\definecolor{furutaLinkFill}{HTML}{8CB4E8}
\definecolor{furutaLinkHi}{HTML}{DCEAFE}
\definecolor{furutaPostEdge}{HTML}{334155}
\definecolor{furutaPostFill}{HTML}{CBD5E1}
\definecolor{furutaPostHi}{HTML}{F1F5F9}
\definecolor{furutaAngle}{HTML}{D9480F}
\definecolor{furutaTorque}{HTML}{0D8A47}
\definecolor{furutaRef}{HTML}{2563EB}
\definecolor{furutaMuted}{HTML}{64748B}
%
\resizebox{0.92\columnwidth}{!}{%
\begin{tikzpicture}[
    >=Stealth,
    line cap=round,
    line join=round,
    font=\small,
    % Estilos semanticos unificados
    axis0/.style={->, >={Stealth[length=4.8pt, width=3.4pt]}, line width=0.95pt, furutaSlate},
    dimline/.style={<->, >={Stealth[length=4.0pt, width=2.8pt]}, line width=0.6pt, furutaSlate!85},
    dimext/.style={line width=0.45pt, furutaSlate!45},
    dimlbl/.style={font=\footnotesize, text=furutaSlate, fill=white, inner sep=1.4pt},
    refline/.style={dashed, line width=0.85pt, furutaRef!85},
    refsub/.style={dashed, line width=0.65pt, furutaMuted!80},
    angarc/.style={->, >={Stealth[length=4.4pt, width=3.1pt]}, line width=1.15pt, furutaAngle},
    anglbl/.style={font=\footnotesize\bfseries, text=furutaAngle, inner sep=1.5pt},
    halo/.style={preaction={draw=white, line width=2.8pt}}
]
    % =========================================================================
    % 1. PARAMETROS CENTRALIZADOS DE PROYECCION 3D Y GEOMETRIA (Single Source)
    % =========================================================================
    % Camara ortografica: azimut phi_c = 28 deg, elevacion psi_c = 22 deg
    \pgfmathsetmacro{\phiCam}{28}
    \pgfmathsetmacro{\psiCam}{22}
    % Brazo horizontal en theta = 60 deg; desviacion invertida tilde_alpha = 33 deg
    % Coherente con Eq. (1): alpha = pi + tilde_alpha (tilde_alpha > 0 -> +j1)
    \pgfmathsetmacro{\thetaArm}{60}
    \pgfmathsetmacro{\alphaTilde}{33}

    % Vectores base proyectados de F0 = {X0, Y0, Z0}
    \pgfmathsetmacro{\Xx}{-sin(\phiCam)}
    \pgfmathsetmacro{\Xy}{-sin(\psiCam)*cos(\phiCam)}
    \pgfmathsetmacro{\Yx}{cos(\phiCam)}
    \pgfmathsetmacro{\Yy}{-sin(\psiCam)*sin(\phiCam)}
    \pgfmathsetmacro{\Zx}{0}
    \pgfmathsetmacro{\Zy}{cos(\psiCam)}

    % Marco local del brazo: i1 (radial) y j1 = k0 x i1 (tangencial +theta)
    \pgfmathsetmacro{\Xr}{cos(\thetaArm)*\Xx + sin(\thetaArm)*\Yx}
    \pgfmathsetmacro{\Yr}{cos(\thetaArm)*\Xy + sin(\thetaArm)*\Yy}
    \pgfmathsetmacro{\Xt}{-sin(\thetaArm)*\Xx + cos(\thetaArm)*\Yx}
    \pgfmathsetmacro{\Yt}{-sin(\thetaArm)*\Xy + cos(\thetaArm)*\Yy}

    % Direccion unitaria del brazo en pantalla y normal perpendicular hacia abajo
    \pgfmathsetmacro{\alen}{sqrt(\Xr*\Xr + \Yr*\Yr)}
    \pgfmathsetmacro{\uaX}{\Xr/\alen}
    \pgfmathsetmacro{\uaY}{\Yr/\alen}
    \pgfmathsetmacro{\naX}{\uaY}
    \pgfmathsetmacro{\naY}{-\uaX}

    % Direccion del pendulo: p(pi + tilde_alpha) = sin(tilde_alpha)*j1 + cos(tilde_alpha)*k0
    \pgfmathsetmacro{\uxP}{sin(\alphaTilde)*\Xt + cos(\alphaTilde)*\Zx}
    \pgfmathsetmacro{\uyP}{sin(\alphaTilde)*\Yt + cos(\alphaTilde)*\Zy}
    \pgfmathsetmacro{\plen}{sqrt(\uxP*\uxP + \uyP*\uyP)}
    \pgfmathsetmacro{\upX}{\uxP/\plen}
    \pgfmathsetmacro{\upY}{\uyP/\plen}
    % Normal unitaria hacia la derecha del pendulo en pantalla
    \pgfmathsetmacro{\npX}{\upY}
    \pgfmathsetmacro{\npY}{-\upX}

    % Longitudes, radios y holguras parametricas (escala ~1:1 en columna IEEE)
    \def\lenX{3.20}
    \def\lenY{2.25}
    \def\lenZ{2.55}
    \def\lenPost{2.35}
    \def\lenArm{6.20}
    \def\lenPend{4.25}
    \def\fracCM{0.56}
    \def\lenDownRef{1.15}

    \def\rTau{0.50}
    \def\zTau{1.28}
    \def\rTheta{1.65}
    \def\rTilde{3.08}
    \def\rCM{0.115}

    \def\dArm{1.32}
    \def\dLone{0.62}
    \def\dLtot{1.16}
    \def\extOver{0.12}

    % =========================================================================
    % 2. UNICO ORIGEN (O0) Y COORDENADAS NOMBRADAS DERIVADAS POR CONSTRUCCION
    % =========================================================================
    \coordinate (O0)       at (0, 0);
    \coordinate (X0tip)    at ($(O0) + ({\lenX*\Xx}, {\lenX*\Xy})$);
    \coordinate (Y0tip)    at ($(O0) + ({\lenY*\Yx}, {\lenY*\Yy})$);
    \coordinate (Z0tip)    at ($(O0) + ({\lenZ*\Zx}, {\lenZ*\Zy})$);
    \coordinate (PostBot)  at ($(O0) - ({\lenPost*\Zx}, {\lenPost*\Zy})$);
    \coordinate (TauCtr)   at ($(O0) + ({\zTau*\Zx}, {\zTau*\Zy})$);

    \coordinate (Ppiv)     at ($(O0) + ({\lenArm*\Xr}, {\lenArm*\Yr})$);
    \coordinate (Pup)      at ($(Ppiv) + ({\lenPend*\Zx}, {\lenPend*\Zy})$);
    \coordinate (Pdown)    at ($(Ppiv) - (0, \lenDownRef)$);
    \coordinate (Pcm)      at ($(Ppiv) + ({\fracCM*\lenPend*\uxP}, {\fracCM*\lenPend*\uyP})$);
    \coordinate (Ptip)     at ($(Ppiv) + ({\lenPend*\uxP}, {\lenPend*\uyP})$);

    % Puntos de cota para L0 (perpendiculares al brazo en pantalla)
    \coordinate (L0ExtO)   at ($(O0)   + ({(\dArm+\extOver)*\naX}, {(\dArm+\extOver)*\naY})$);
    \coordinate (L0ExtP)   at ($(Ppiv) + ({(\dArm+\extOver)*\naX}, {(\dArm+\extOver)*\naY})$);
    \coordinate (L0DimO)   at ($(O0)   + ({\dArm*\naX}, {\dArm*\naY})$);
    \coordinate (L0DimP)   at ($(Ppiv) + ({\dArm*\naX}, {\dArm*\naY})$);

    % Puntos de cota escalonadas para l1 y L1 (lado derecho del pendulo)
    \coordinate (PivExtL1) at ($(Ppiv) + ({(\dLtot+\extOver)*\npX}, {(\dLtot+\extOver)*\npY})$);
    \coordinate (CmExt)    at ($(Pcm)  + ({(\dLone+\extOver)*\npX}, {(\dLone+\extOver)*\npY})$);
    \coordinate (TipExt)   at ($(Ptip) + ({(\dLtot+\extOver)*\npX}, {(\dLtot+\extOver)*\npY})$);
    \coordinate (l1DimA)   at ($(Ppiv) + ({\dLone*\npX}, {\dLone*\npY})$);
    \coordinate (l1DimB)   at ($(Pcm)  + ({\dLone*\npX}, {\dLone*\npY})$);
    \coordinate (L1DimA)   at ($(Ppiv) + ({\dLtot*\npX}, {\dLtot*\npY})$);
    \coordinate (L1DimB)   at ($(Ptip) + ({\dLtot*\npX}, {\dLtot*\npY})$);

    % =========================================================================
    % 3. LINEAS DE EXTENSION, COTA L0 (DETRAS DEL POSTE) Y REFERENCIAS
    % =========================================================================
    \draw[dimext] (O0)   -- (L0ExtO);
    \draw[dimext] (Ppiv) -- (L0ExtP);
    \draw[dimext] (Ppiv) -- (PivExtL1);
    \draw[dimext] (Pcm)  -- (CmExt);
    \draw[dimext] (Ptip) -- (TipExt);

    % Cota paralela limpia L0 bajo el brazo (pasa limpiamente detras del poste -Z0)
    \draw[dimline] (L0DimO) -- node[dimlbl, midway, sloped] {$L_0$} (L0DimP);

    % Referencia sutil inferior (alpha = 0)
    \draw[refsub] (Ppiv) -- (Pdown);
    \node[anchor=north, font=\scriptsize, text=furutaMuted, inner sep=1.5pt]
        at ($(Pdown) + (0, -0.04)$) {$\alpha = 0$};

    % Mitad trasera del arco de par tau_0 (detras del eje +Z0)
    \begin{scope}[shift={(TauCtr)}]
        \draw[line width=1.45pt, furutaTorque]
            plot[domain=62:201, samples=26, smooth]
            ({ \rTau*(cos(\x)*\Xx + sin(\x)*\Yx) },
             { \rTau*(cos(\x)*\Xy + sin(\x)*\Yy) });
    \end{scope}

    % =========================================================================
    % 4. POSTE VERTICAL (-Z0) Y EJES ORTOGONALES INERCIALES F0
    % =========================================================================
    % Poste cilindrico vertical del soporte motor (sin dibujar motor ni caja)
    \draw[white,          line width=6.2pt, line cap=round] (O0) -- (PostBot);
    \draw[furutaPostEdge, line width=4.6pt, line cap=round] (O0) -- (PostBot);
    \draw[furutaPostFill, line width=3.0pt, line cap=round] (O0) -- (PostBot);
    \draw[furutaPostHi,   line width=0.9pt, line cap=round] ($(O0)+(-0.025,0)$) -- ($(PostBot)+(-0.025,0)$);

    % Ejes inerciales ortogonales (+X0, +Y0, +Z0) con halo sutil
    \draw[axis0, halo] (O0) -- (X0tip)
        node[below left=-1pt and 1pt, font=\small\bfseries, text=furutaSlate] {$+X_0$};
    \draw[axis0, halo] (O0) -- (Y0tip)
        node[right=2pt, font=\small\bfseries, text=furutaSlate] {$+Y_0$};
    \draw[axis0, halo] (O0) -- (Z0tip)
        node[above=1pt, font=\small\bfseries, text=furutaSlate] {$+Z_0$};

    % Mitad frontal del arco de par tau_0 (emerge a la izquierda tras +Z0 sin muesca)
    \begin{scope}[shift={(TauCtr)}]
        \draw[->, >={Stealth[length=4.6pt, width=3.2pt]}, line width=1.45pt, furutaTorque]
            plot[domain=215:336, samples=26, smooth]
            ({ \rTau*(cos(\x)*\Xx + sin(\x)*\Yx) },
             { \rTau*(cos(\x)*\Xy + sin(\x)*\Yy) });
    \end{scope}
    \node[anchor=south east, font=\small\bfseries, text=furutaTorque, inner sep=1.5pt]
        at ($(TauCtr) + ({-\rTau*0.62}, {\rTau*0.58})$) {$\tau_0$};

    % Etiqueta limpia del origen O0 en cuadrante superior izquierdo despejado
    \node[anchor=south east, font=\footnotesize\bfseries, text=furutaSlate, inner sep=1.5pt]
        at ($(O0) + (-0.12, 0.08)$) {$O_0$};

    % =========================================================================
    % 5. ARCO ANGULAR theta(t) Y BRAZO HORIZONTAL (L0)
    % =========================================================================
    % Arco angular theta(t) desde +X0 hasta el brazo horizontal en plano X0-Y0
    \begin{scope}[shift={(O0)}]
        \draw[angarc, halo]
            plot[domain=4:{\thetaArm-5.5}, samples=30, smooth]
            ({ \rTheta*(cos(\x)*\Xx + sin(\x)*\Yx) },
             { \rTheta*(cos(\x)*\Xy + sin(\x)*\Yy) });
    \end{scope}
    \pgfmathsetmacro{\thetaLbl}{0.56*\thetaArm}
    \node[anglbl, anchor=north, fill=white, inner sep=1pt]
        at ($(O0) + ({(\rTheta+0.12)*(cos(\thetaLbl)*\Xx + sin(\thetaLbl)*\Yx)},
                     {(\rTheta+0.12)*(cos(\thetaLbl)*\Xy + sin(\thetaLbl)*\Yy)})$)
        {$\theta(t)$};

    % Varilla cilindrica esbelta del brazo horizontal (L0)
    \draw[white,          line width=6.2pt, line cap=round] (O0) -- (Ppiv);
    \draw[furutaLinkEdge, line width=4.4pt, line cap=round] (O0) -- (Ppiv);
    \draw[furutaLinkFill, line width=2.8pt, line cap=round] (O0) -- (Ppiv);
    \draw[furutaLinkHi,   line width=0.85pt, line cap=round]
        ($(O0)+(0.015, 0.025)$) -- ($(Ppiv)+(0.015, 0.025)$);

    % Junta esferica/cilindrica limpia en el origen O0
    \fill[furutaLinkFill, draw=furutaLinkEdge, line width=0.85pt] (O0) circle (0.11cm);
    \fill[furutaLinkHi] ($(O0)+(-0.03, 0.03)$) circle (0.032cm);

    % =========================================================================
    % 6. ARTICULACION DISTAL, VERTICAL (alpha = pi), PENDULO (L1) Y tilde_alpha
    % =========================================================================
    % Linea punteada vertical superior de equilibrio invertido (alpha = pi)
    \draw[refline] (Ppiv) -- (Pup);
    \node[anchor=south, font=\footnotesize\bfseries, text=furutaRef!90!black, inner sep=1.5pt]
        at ($(Pup) + (0, 0.06)$) {$\alpha = \pi$};

    % Arco angular limpio tilde_alpha entre la vertical superior y el pendulo
    \begin{scope}[shift={(Ppiv)}]
        \draw[angarc]
            plot[domain=0.4:{\alphaTilde-1.8}, samples=24, smooth]
            ({ \rTilde*(sin(\x)*\Xt + cos(\x)*\Zx) },
             { \rTilde*(sin(\x)*\Yt + cos(\x)*\Zy) });
    \end{scope}
    \pgfmathsetmacro{\midTilde}{0.50*\alphaTilde}
    \node[anglbl, anchor=south, inner sep=1.5pt]
        at ($(Ppiv) + ({(\rTilde+0.12)*(sin(\midTilde)*\Xt + cos(\midTilde)*\Zx)},
                       {(\rTilde+0.12)*(sin(\midTilde)*\Yt + cos(\midTilde)*\Zy)})$)
        {$\tilde{\alpha}$};

    % Varilla cilindrica esbelta del pendulo (L1)
    \draw[white,          line width=6.2pt, line cap=round] (Ppiv) -- (Ptip);
    \draw[furutaLinkEdge, line width=4.4pt, line cap=round] (Ppiv) -- (Ptip);
    \draw[furutaLinkFill, line width=2.8pt, line cap=round] (Ppiv) -- (Ptip);
    \draw[furutaLinkHi,   line width=0.85pt, line cap=round]
        ($(Ppiv)+(-0.025, 0.015)$) -- ($(Ptip)+(-0.025, 0.015)$);

    % Junta circular en articulacion distal Ppiv
    \fill[furutaLinkFill, draw=furutaLinkEdge, line width=0.85pt] (Ppiv) circle (0.11cm);
    \fill[furutaLinkHi] ($(Ppiv)+(-0.03, 0.03)$) circle (0.032cm);

    % Marcador limpio de centro de masa G1 sobre la varilla a distancia l1
    \fill[white, draw=furutaSlate, line width=0.75pt] (Pcm) circle (\rCM cm);
    \fill[furutaSlate] (Pcm) -- ++(0:\rCM cm) arc (0:90:\rCM cm) -- cycle;
    \fill[furutaSlate] (Pcm) -- ++(180:\rCM cm) arc (180:270:\rCM cm) -- cycle;
    \node[anchor=south east, font=\footnotesize\bfseries, text=furutaSlate, inner sep=1.5pt]
        at ($(Pcm) + (-0.11, 0.05)$) {$G_1$};

    % Cotas paralelas escalonadas del pendulo: l1 (interna) y L1 (externa)
    \draw[dimline] (l1DimA) -- node[dimlbl, midway, sloped] {$l_1$} (l1DimB);
    \draw[dimline] (L1DimA) -- node[dimlbl, midway, sloped] {$L_1$} (L1DimB);

    % =========================================================================
    % 7. VECTOR DE GRAVEDAD g EN ESQUINA SUPERIOR DERECHA DESPEJADA
    % =========================================================================
    \node[fit=(Ptip)(TipExt)(Pup), inner sep=0pt] (TopRightScope) {};
    \coordinate (Gtop) at ($(TopRightScope.east |- Pup) + (0.28, -0.05)$);
    \coordinate (Gbot) at ($(Gtop) + (0, -0.95)$);
    \draw[->, >={Stealth[length=5.0pt, width=3.5pt]}, line width=1.15pt, furutaSlate]
        (Gtop) -- (Gbot);
    \node[anchor=west, font=\small\bfseries, text=furutaSlate, inner sep=2pt]
        at ($(Gtop)!0.48!(Gbot) + (0.04, 0)$) {$\vec{g}$};

\end{tikzpicture}%
}%
\caption{Diagrama cinem\'atico y marcos de referencia del p\'endulo de Furuta.}
\label{fig:cinematica_furuta}
\end{figure}
